% Nouveau EGA V (cote 25): archive cote 25, site page 173 (PDF page 174).
% typescript. Draft compared: none.
% AI-transcribed by Claude Opus 5.5 (claude-opus-5-5, adaptive thinking, effort high), Message Batches API, batch ega25b (msgbatch_01A7n5YuXubUPNK1zsepi8RG), 2026-11-13; not yet checked by eye.
\scanpage{25/173}
% NOTE: Both scan halves contain a very faint forward-reading impression. Its physical origin cannot be established from these images. The readings below follow this impression and remain explicitly doubtful; neighbouring source text supplied comparison proposals only.
\var{\margin{\unsure{\hand{IV|//}}}}{\placedmargin{\unsure{\hand{IV|//}}}}{marginal note placed}{ega25_25_173-v0001} \unsure{\hand{nouveau EGA V}}
\unsure{17.16. \uline{Formes quadratiques lisses}.}
% NOTE: The upper-right area yields no secure authorial folio; the irregular outer-edge traces have no securely readable additional text.
\unsure{16.1.Soit $Q$ une section de $\underline{\mathrm{Sym}}^2(\underline{E})$, ($\underline{E}$ loc libre de rang fini),}
\unsure{de sorte que Q \var{\xout{\unsure{définit}}}{}{typewriter strike-out}{ega25_25_173-v0002} (qu'on peut aussi interpréter comme une forme}
\unsure{quadratique sur $\overset{\hand{\vee}}{\underline{E}}$ ) \var{\xout{\illegible{3} \unsure{section}}}{}{typewriter strike-out}{ega25_25_173-v0003} définit une forme biliné-}
\unsure{aire symétrique B sur $\overset{\hand{\vee}}{\underline{E}}$, d'où un homomorphisme $\overset{\hand{\vee}}{\underline{E}} \hand{\longrightarrow} \underline{E}$ , et passant aux}
\unsure{Modules déterminants on trouve $\text{dét}\ \overset{\hand{\vee}}{\underline{E}} \hand{\longrightarrow} \det \underline{E}$ , d'où enfin une section}
\[ \unsure{d(Q) \hand{\in} \hand{\Gamma}(\ \text{dét}\ (E)^{\hand{\otimes} 2})} \]
\unsure{appelée (sauf erreur) \uline{discriminant de la forme quadratique} Q. Lorsque}
\unsure{$\underline{E} = \underline{O}_S^{\ \hand{n}}$ , c'est simplement la section de $\underline{O}_S$ déterminant de la matrice}
\unsure{à n lignes et n colonnes, à coeff dans $\hand{\Gamma}\underline{O}_S$ , qui exprime Q . Notons tout}
\unsure{de suite:}
\unsure{\uline{Proposition} 16.2. Supposons que pour tout $s \hand{\in} S$, on ait car $k(s) \hand{\neq} 2$ \uline{ou}}
\unsure{rang. $E\hand{\otimes} k(s)$ pair. Alors, pour que Q soit lisse, il f et s que Q soit}
\unsure{``non dégénérée'' i.e. que d(Q) soit inversible. \hand{(}D'ailleurs la condition est}
\unsure{suffisante sans restriction sur S ni $\underline{E}$\hand{)}.}
\unsure{La dernière assertion n'est mise que pour la commodité des référen-}
\unsure{ces éventuelles, car elle est trivialement contenue dans la première,}
\unsure{l'hypothèse que Q est non dégénérée implique en effet que $\underline{E}$ est de rang}
\unsure{pair en tout point dont la car est 2. D'ailleurs, elle se démontre direc-}
\unsure{tement sans considération de caractéristique, car (se plaçant sur un corps}
\unsure{\var{\inserted{\unsure{alg clos}}}{\unsure{alg clos}}{insertion}{ega25_25_173-v0004} et prenant une base) elle exprime que les $\frac{\hand{\partial} Q}{\hand{\partial} X_i}$ n'ont pas de zéro\var{\xout{\unsure{s}}}{}{typewriter strike-out}{ega25_25_173-v0005} commun\var{\xout{\unsure{s}}}{}{typewriter strike-out}{ega25_25_173-v0006}}
\unsure{non trivial, à fortiori \unsure{ils} n'ont pas de zéro commun non trivial avec Q.}
\unsure{Si car $k \hand{\neq} 2$, la réciproque est vraie, car en vertu de la formule}
\[ \unsure{2Q = \hand{\sum} X_i \frac{\hand{\partial} Q}{\hand{\partial} X_i}} \]
\unsure{on voit alors que tout zéro commun aux $\frac{\hand{\partial} \unsure{Q|E}}{\hand{\partial} X_i}$ est aussi un zéro de Q. Enfin,}
\unsure{si car $k = 2$ , \var{\xout{\unsure{si E est de rang pair}}}{}{typewriter strike-out}{ega25_25_173-v0007} alors la forme bilinéaire B}
\mspage{\hand{173}}
